\chapter*{FORMAT DEKLARASI PENGGUNAAN AI}
\addcontentsline{toc}{chapter}{FORMAT DEKLARASI PENGGUNAAN AI}

\begingroup\setlength{\cmidrulekern}{1.5em}\singlespacing\setlength{\parskip}{3pt}
\vspace{-2.4em}
{\centering\bfseries SURAT PERNYATAAN PENGGUNAAN AI\par}

Dengan ini, kami:

{\renewcommand{\arraystretch}{1.3}\hspace*{1.5em}\begin{tabular}{@{}l@{\hspace{0.5em}}c@{\hspace{0.5em}}l@{}}
Nama/NIM & : & \underline{\makebox[0.68\linewidth][l]{Shanice Feodora Tjahjono / 13523097}} \\
 &  & \underline{\makebox[0.68\linewidth][l]{Brian Ricardo Tamin / 13523126}} \\
 &  & \underline{\makebox[0.68\linewidth][l]{Nathanael Rachmat / 13523142}} \\
Fakultas/Sekolah & : & \underline{\makebox[0.68\linewidth][l]{Sekolah Teknik Elektro dan Informatika}} \\
Kode Mata Kuliah & : & \underline{\makebox[0.68\linewidth][l]{IF4031}} \\
Nama Mata Kuliah & : & \underline{\makebox[0.68\linewidth][l]{Arsitektur Aplikasi Terdistribusi}} \\
Judul Tugas/Proposal & : & \underline{\makebox[0.68\linewidth][l]{Platform Koordinasi Kebencanaan (Tugas Besar Milestone 1, Kelompok Ingress)}} \\
\end{tabular}}

Menyatakan bahwa kami menggunakan AI dalam pengerjaan maupun penyusunan luaran tugas pada mata kuliah yang tertulis di atas.

\vspace{\baselineskip}
Jika Ya, maka alat AI/kecerdasan buatan yang digunakan adalah:

{\setlength{\LTpre}{0pt}\setlength{\LTpost}{0pt}\setlength{\LTleft}{0pt}\setlength{\LTright}{0pt}\small\renewcommand{\arraystretch}{1}\setlength{\extrarowheight}{6pt}\setlength{\tabcolsep}{6pt}
\begin{longtable}[]{|>{\raggedright\arraybackslash}m{0.58\linewidth}|>{\centering\arraybackslash}m{0.15\linewidth}|>{\centering\arraybackslash}m{0.15\linewidth}|}
\hline
\multicolumn{1}{|c|}{\textbf{Lingkup Pekerjaan}} & \multicolumn{1}{c|}{\textbf{Digunakan?}} & \multicolumn{1}{c|}{\textbf{Tingkat Penggunaan}} \\[6pt]
\hline
\endhead
\rule{0pt}{12pt}\textbf{Pemeriksaan Ejaan:} menggunakan \textit{tools} seperti Grammarly, Paperpal, Deepl, Quillbot, Grammarbot, LanguageTool, ProWritingAid, ChatGPT, Google Gemini, atau sejenisnya.\rule[-8pt]{0pt}{0pt} & \rule{0pt}{12pt}Tidak\rule[-8pt]{0pt}{0pt} & \rule{0pt}{12pt}1\rule[-8pt]{0pt}{0pt} \\
\hline
\rule{0pt}{12pt}\textbf{Pembuatan Teks:} menggunakan \textit{tools} seperti ChatGPT, Gemini, Paperpal, Copilot, GrammarlyGO, WordAI, WriteSonic, Jasper, Jenni AI, atau sejenisnya.\rule[-8pt]{0pt}{0pt} & \rule{0pt}{12pt}Ya\rule[-8pt]{0pt}{0pt} & \rule{0pt}{12pt}4\rule[-8pt]{0pt}{0pt} \\
\hline
\rule{0pt}{12pt}\textbf{Bantuan Diskusi dalam Penyusunan Konten:} menggunakan \textit{tools} seperti ChatGPT, Gemini, Copilot, Perplexity, Jenni AI, Zoom-Companion, atau sejenisnya.\rule[-8pt]{0pt}{0pt} & \rule{0pt}{12pt}Ya\rule[-8pt]{0pt}{0pt} & \rule{0pt}{12pt}2\rule[-8pt]{0pt}{0pt} \\
\hline
\rule{0pt}{12pt}\textbf{Pembuatan Gambar, Video, dan/atau Grafik:} menggunakan \textit{tools} seperti Craiyon, DALL-E, Midjourney, Stable Diffusion, Microsoft Designer, Gemini, Canva AI, atau sejenisnya.\rule[-8pt]{0pt}{0pt} & \rule{0pt}{12pt}Tidak\rule[-8pt]{0pt}{0pt} & \rule{0pt}{12pt}1\rule[-8pt]{0pt}{0pt} \\
\hline
\rule{0pt}{12pt}\textbf{Lainnya, Jelaskan:} Claude Code juga dipakai untuk alur kerja repositori, yaitu mengelola \textit{issue}, \textit{pull request}, dan graf dependensi \textit{backlog} lewat \texttt{gh}, serta merapikan komentar kode dan menyusun bab laporan. Aturan kerjanya tertulis pada \texttt{AGENTS.md}, \texttt{doc/latex/AGENTS.md}, dan \textit{skill} Go serta \textit{antislop} di \texttt{.claude/skills}.\rule[-8pt]{0pt}{0pt} & \rule{0pt}{12pt}Ya\rule[-8pt]{0pt}{0pt} & \rule{0pt}{12pt}3\rule[-8pt]{0pt}{0pt} \\
\hline
\end{longtable}}
\vspace{-0.2em}

Kami juga menyatakan bahwa setiap penggunaan AI/kecerdasan buatan yang dilakukan pada mata kuliah tersebut telah dinyatakan di dalam surat ini.


\vspace{\baselineskip}
\hfill Bandung, 9 Oktober 2026\hspace*{1.5em}

\hfill\begin{tabular}{@{}>{\raggedright\arraybackslash}p{0.27\linewidth}@{\hspace{1em}}>{\raggedright\arraybackslash}p{0.27\linewidth}@{\hspace{1em}}>{\raggedright\arraybackslash}p{0.27\linewidth}@{}}
\rule{0pt}{1.8cm} & & \\
\cmidrule(r){1-1}\cmidrule(r){2-2}\cmidrule(r){3-3}
Shanice Feodora Tjahjono & Brian Ricardo Tamin & Nathanael Rachmat \\
NIM 13523097 & NIM 13523126 & NIM 13523142 \\
\end{tabular}
\endgroup
